% ===================================================================
%  جدول نمبر 11 — شہر اور اُس کی عمارتیں۔ --- آتش زدگی۔
%
%  Traduction, faite en 2026, en ourdou d'aujourd'hui, sur l'ido de
%  texto/io. Regles de la colonne : voir 10-jadval-01.tex. Deux
%  parties, deux numerotations : les cles portent c1 et c2. Ecrit
%  avec outils/ordo.py sous les yeux.
%
%  LE TABLEAU DE L'ADMINISTRATION EST CELUI OU L'OURDOU N'A RIEN A
%  INVENTER, ET LA RAISON EST HISTORIQUE, PAS LINGUISTIQUE. Le
%  prefet de l'ido est un administrateur de departement nomme par
%  l'Etat ; l'Inde britannique a bati exactement le meme office sur
%  le meme modele, et il porte en ourdou un nom que tout le monde
%  connait : ڈپٹی کمشنر. La colonne l'ecrit sans detour, et ecrit de
%  meme محصول خانہ la douane, بلدیہ کا گھر la mairie, عدالت le
%  tribunal, بچت خانہ la caisse d'epargne, ڈاک خانہ la poste,
%  کتب خانہ la bibliotheque, عجائب گھر le musee, دار المعلمین
%  l'ecole normale, گھنٹہ گھر le beffroi. Aucun de ces mots n'est un
%  calque de circonstance : ce sont les enseignes qu'on lit encore
%  dans les villes du Pendjab.
%
%  گھنٹہ گھر EST LE PLUS EXACT DES NEUF. Le beffroi de l'ido est une
%  tour communale qui porte l'horloge et la cloche et qui domine la
%  vieille ville. Le گھنٹہ گھر de Faisalabad, de Multan, de Lahore
%  est cela meme, bati au meme siecle et pour la meme raison — et le
%  mot dit litteralement « la maison de la cloche ». C'est la
%  premiere fois de tout le releve qu'un monument du fac-simile
%  trouve dans une langue non europeenne un equivalent qui soit a la
%  fois le mot juste, la chose juste et l'epoque juste.
%
%  LE VIOLON (81) NE DEVIENT PAS UNE SARANGI. سارنگی est l'instrument
%  a archet du nord de l'Inde, il est plus proche du violon que rien
%  d'autre, et la tentation etait forte puisque le mot est vivant la
%  ou وائلن est un emprunt. Mais la sarangi n'a ni le meme nombre de
%  cordes, ni le meme jeu, ni la meme place : elle n'est pas dans un
%  orchestre de fosse. La colonne ecrit وائلن, چیلو, ٹرومبون. Elle
%  garde en revanche بانسری pour la flute (82) et ڈھول pour le
%  tambour (85), parce que ceux-la SONT les choses. Neuvieme refus de
%  ce genre, et le premier ou le mot refuse etait le plus beau.
%
%  ET LES POMPIERS N'ONT PAS DE MOT COURT. L'ourdou dit
%  آگ بجھانے والا, « celui qui eteint le feu » — quatre syllabes de
%  periphrase la ou l'ido a « pumpisto ». Ce n'est pas un trou de
%  vocabulaire mais un trou d'institution : le corps de pompiers
%  municipal est arrive tard dans les villes indiennes, et
%  فائر بریگیڈ, l'emprunt anglais qui sert aujourd'hui, nomme le
%  service, pas l'homme. La colonne prend la periphrase plutot que
%  l'emprunt : elle est plus longue mais elle se lit.
%
%  ET LA MEME FAUTE QU'AU TABLEAU 9, DEUX TABLEAUX PLUS LOIN : une
%  cle inventee, t11-c1-12-2, la ou l'ido remet un \VUblancAlinea au
%  milieu du bloc t11-c1-12-1 parce qu'une fin de page tombe la. Deux
%  alineas, une seule cle. Cette fois la faute a servi : cinq blocs
%  du volume sont dans ce cas — t01-09-2, t07-c1-07-1, t09-c1-05-1,
%  t11-c1-12-1, t11-c2-03-1 — et ordo.py les marque desormais d'un
%  « ¶2 ». Un outil qui prevoit vaut mieux qu'un controle qui punit :
%  renvoji.py savait dire « cle inconnue », il ne savait pas dire
%  « n'en ouvre pas une seconde ».
%
%  VINGT ET UNE PAIRES, ET AUCUNE FAUTE D'ORDRE AU PREMIER JET.
%  Deuxieme tableau de suite : la relative postposee a servi quatre
%  fois — la porte (8) avant le chateau (7), la coupole (40) avant
%  l'hopital (39), la poste (43) avant la maison privee (42), les
%  musiciens (81) avant l'orchestre (80). La regle annoncee au
%  tableau 9 se confirme deux fois ; il reste a la voir tenir sur un
%  tableau qui DECRIT au lieu d'enumerer, comme le 9 et le 16.
%
%  Precision, parce que la phrase ci-dessus serait fausse sans elle :
%  renvoji.py a bien rendu DEUX divergences au premier passage, mais
%  les deux venaient de la seule cle inventee dite plus bas, pas d'un
%  renvoi mal place. Zero faute d'ordre n'est pas zero divergence, et
%  on ne compte pas les unes pour les autres.
% ===================================================================

%%K t11-apar-1 apar
\VUsaut{2.0mm}
\VUcentre{13.2pt}{90}{تیسرا سلسلہ}
\VUsaut{3.0mm}
\VUfilet{16mm}
\VUsaut{5.6mm}
\VUtitre{15.0pt}{60}{جدول نمبر 11}
\VUsaut{1.4mm}
\VUpk{11.4pt}{40}{شہر اور اُس کی عمارتیں۔ --- آتش زدگی۔}
\VUsaut{2.2mm}

%%K t11-tit-1 sub
\VUpk{10.2pt}{40}{مضافات۔}
\VUsaut{3.0mm}

%%K t11-c1-01-1 p
1. --- میں ایک بڑے شہر میں رہتا ہوں جو سمندر سے سو کلومیٹر کے فاصلے
پر ہے، اور پھر بھی اونچے درجے کی \VUgras{بندرگاہ}
\textsuperscript{(1)} ہے۔ جو دریا اُس کے کنارے بہتا ہے، وہ چار سو میٹر
چوڑا ہے اور ایک نہایت خوبصورت موڑ بناتا ہے۔

%%K t11-c1-02-1 p
2. --- شہر کا سارا وہ حصہ جو \VUgras{گھاٹوں} \textsuperscript{(2)} پر
ہے، بالکل سمندری اور صنعتی دکھائی دیتا ہے، اور ایک ایسا
\VUgras{مضافاتی محلہ} \textsuperscript{(3)} بناتا ہے جس میں صرف ملّاح
اور مزدور رہتے ہیں۔ یہ لوگ \VUgras{کارخانوں} \textsuperscript{(4)} میں
اور \VUgras{گیس گھر} \textsuperscript{(5)} میں کام کرتے ہیں، جس کی
اونچی \VUgras{چمنی} \textsuperscript{(6)} اور جس کا \VUgras{گیس کا
حوض} بہت دور سے دکھائی دیتے ہیں۔

%%K t11-c1-03-1 p
3. --- قرونِ وسطیٰ میں شہر قلعہ بند تھا، اور اُس کی فصیلوں کے کچھ
باقیات آج بھی ملتے ہیں — خاص طور پر وہ \VUgras{دروازہ}
\textsuperscript{(8)} جو پرانے \VUgras{مضبوط قلعے} \textsuperscript{(7)}
کا ہے اور جس کے دونوں طرف اُس کے دو \VUgras{بُرج}
\textsuperscript{(9)} کھڑے ہیں ؛ وہ برج اب \VUgras{قید خانے} کے طور پر
کام آتے ہیں۔

%%K t11-c1-04-1 p
4. --- اِس سارے \VUgras{محلے} میں، جو بہت پرانا لگتا ہے، صرف ایک
عمارت نسبتاً جدید ہے : \VUgras{محصول خانہ} \textsuperscript{(10)}۔ وہ
گھاٹ اور اُس \VUgras{گلی} \textsuperscript{(11)} کے بیچ واقع ہے جو
پرانے \VUgras{بلدیہ کے گھر} \textsuperscript{(12)} کو جاتی ہے۔ اِس
آخری عمارت کو اُس دیو قامت \VUgras{گھنٹہ گھر} \textsuperscript{(13)}
سے آسانی سے پہچان لیا جاتا ہے جو قدیم شہر پر حاوی ہے۔

%%K t11-c1-05-1 p
5. --- سڑک کی دوسری طرف ایک عمارت \VUgras{کھڑی} ہے، جو محصول خانے ہی
کے انداز میں بنی ہے : یہ \VUgras{اسٹاک ایکسچینج} \textsuperscript{(14)}
ہے۔ اُس کا ایک رُخ ایک چھوٹے چوک کی طرف دیکھتا ہے، جہاں
\VUgras{ڈپٹی کمشنر کا دفتر} \textsuperscript{(15)} ہے۔ حقیقت یہ ہے کہ
ہمارا شہر دار الحکومت نہ ہونے کے باوجود ضلعے کا اہم صدر مقام ہے۔

%%K t11-tit-2 sub
\VUsaut{5.0mm}
\VUpk{10.2pt}{40}{شہر کا سب سے اونچا حصہ۔}
\VUsaut{3.4mm}

%%K t11-c1-06-1 p
6. --- چھاؤنی، جو خاصی بڑی ہے، بہت صحت بخش اور کشادہ
\VUgras{بیرکوں} \textsuperscript{(16)} میں ٹھہری ہوئی ہے ؛ وہ بیرکیں
\VUgras{بلدیاتی باغ} \textsuperscript{(17)} کی جالی تک پھیلی ہوئی ہیں۔
جو \VUgras{شاہراہ} \textsuperscript{(19)} اِس باغ کو جاتی ہے، وہ
\VUgras{بچت خانے} \textsuperscript{(18)} کے پیچھے سے گزرتی ہے اور
\VUgras{بڑے کلیسا} \textsuperscript{(20)} کے چوک پر جا نکلتی ہے۔

%%K t11-c1-07-1 p
7. --- یہ ساری عمارتیں ایک ٹیلے پر حلقہ در حلقہ پھیلی ہوئی ہیں، جہاں
ہمارا کشادہ \VUgras{ثانوی اسکول} \textsuperscript{(21)} بھی ہے۔

%%K t11-c1-07-2 p
اگر کوئی ذرا ڈھلوان راستے سے اترے، جو بڑی سڑک سے جا ملتا ہے، تو
\VUgras{عدالت} \textsuperscript{(22)} تک پہنچ جاتا ہے، جس کے سامنے ایک
خوشگوار \VUgras{عوامی باغ} \textsuperscript{(26)} پھیلا ہوا ہے۔ یہ
\VUgras{باغیچہ} بہت بڑا نہیں، مگر شہر کے بیچ میں ہریالی اور ٹھنڈک کا
ایک نخلستان بنا دیتا ہے۔ اِسی لیے شہر والوں کا وہاں بہت آنا جانا ہے ؛
وہ \VUgras{قہوہ خانے} \textsuperscript{(25)} کے شامیانے کے نیچے آ
بیٹھنا پسند کرتے ہیں تاکہ اُن فوجی ساز نوازوں کو سن سکیں جو جمعرات
اور اتوار کو اُس \VUgras{چھتری} \textsuperscript{(27)} میں بجاتے ہیں جو
سبزہ زاروں اور پھولوں کے تختوں کے بیچ رکھی ہے۔

%%K t11-c1-08-1 p
8. --- اُنھی سبزہ زاروں میں سے ایک پر کانسی کا ایک \VUgras{مجسمہ}
\textsuperscript{(28)} کھڑا ہے — ہمارے ایک ہم شہری کی یاد میں کھڑی کی
ہوئی یادگار، جس نے اپنے جنم شہر کی بڑی خدمت کی تھی۔

%%K t11-tit-3 sub
\VUsaut{4.0mm}
\VUpk{10.2pt}{40}{سواریاں۔}
\VUsaut{3.4mm}

%%K t11-c1-09-1 p
9. --- اب تک ہمارا شہر بجلی کی روشنی سے روشن نہیں ؛ اُس کے پاس صرف
\VUgras{گیس کے لیمپ} \textsuperscript{(29)} ہیں۔ مگر \VUgras{یہاں کے
اخبارات} اِس بات کی مہم چلا رہے ہیں کہ روشنی کا یہ نظام بہتر کیا
جائے، \VUgras{جیسے} \VUgras{ٹراموں کے چلانے کا نظام} پہلے ہی بہتر کیا
جا چکا ہے۔ \VUgras{وہ} پرانی گاڑیاں، جنھیں \VUgras{گھوڑے}
\VUgras{کھینچتے} تھے، عملاً \VUgras{بجلی کی ٹراموں}
\textsuperscript{(31)} سے بدل دی گئیں ؛ وہ مسافروں کو \VUgras{شہر کے
ایک سرے سے دوسرے سرے تک} بہت \VUgras{سستے} داموں تیزی سے پہنچاتی ہیں۔

%%K t11-c1-10-1 p
10. --- عدالت گھر کے پاس \VUgras{بینک} \textsuperscript{(32)} ہے، جہاں
تجارتی ہنڈیاں بھنائی جاتی ہیں۔ \VUgras{بینک کے وصولی کار}
\textsuperscript{(33)} گھر گھر جا کر واجبات وصول کرتے ہیں۔

%%K t11-tit-4 sub
\VUsaut{4.0mm}
\VUpk{10.2pt}{40}{فن اور تعلیم۔}
\VUsaut{3.4mm}

%%K t11-c1-11-1 p
11. --- جو خوبصورت عمارت پاس ہی کھڑی ہے، وہ \VUgras{عجائب گھر} ہے۔
اُس میں شہر کا \VUgras{کتب خانہ} \textsuperscript{(34)}،
\VUgras{قدرتی علوم کے} خوبصورت \VUgras{ذخیرے} \textsuperscript{(35)}
(قدرت گھر) اور ایک خوش نما \VUgras{تصویر گھر} \textsuperscript{(36)}
اکٹھے ہیں ؛ یہ آخری ایک شان دار مستقل \VUgras{نمائش گاہ} بناتا ہے۔

%%K t11-c1-11-2 p
وہ ہفتے میں دو بار عوام کے لیے کھلتا ہے، مگر باہر کے لوگ کسی بھی دن
\VUgras{چوکیدار} \textsuperscript{(37)} کو تھوڑا سا انعام دے کر دیکھ
سکتے ہیں۔ \VUgras{مصور} \textsuperscript{(38)} وہاں کام کرنے آتے ہیں۔
اُنھیں اپنی تپائی کے \VUgras{سامنے} دیکھا جا سکتا ہے، ہاتھ میں
\VUgras{رنگ تختی} لیے، مشہور تصویروں کی نقل اتارتے ہوئے۔

%%K t11-c1-12-1 p
12. --- عجائب گھر کے پیچھے، اونچائی پر، وہ \VUgras{گنبد}
\textsuperscript{(40)} دکھائی دینے لگتا ہے جو \VUgras{اسپتال}
\textsuperscript{(39)} کا ہے، اور \VUgras{دار المعلمین}
\textsuperscript{(41)} کی عمارتیں بھی، جس میں ہمارے بلدیاتی اسکولوں کے
اساتذہ کو تربیت دی جاتی تھی۔

تماشا گھر کے چوک پر \VUgras{ڈاک خانہ} \textsuperscript{(43)} ہے، جو
ایک \VUgras{نجی مکان} \textsuperscript{(42)} میں قائم کیا گیا ہے۔

%%K t11-tit-5 sub
\VUsaut{5.0mm}
\VUpk{11.4pt}{40}{آتش زدگی۔}
\VUsaut{3.0mm}

%%K t11-tit-6 sub
\VUpk{10.2pt}{40}{آگ کا شور۔}
\VUsaut{3.4mm}

%%K t11-c2-01-1 p
1. --- ایک دہشت ناک خبر سارے شہر میں پھیل جاتی ہے : تماشا گھر میں
ابھی ابھی آگ لگی ہے! جس لمحے میں \VUgras{چوک} \textsuperscript{(96)}
پر پہنچتا ہوں، مجمع پہلے ہی \VUgras{فٹ پاتھ} \textsuperscript{(46)}
اور \VUgras{پختہ سڑک} \textsuperscript{(47)} پر جمع ہو چکا ہے۔
\VUgras{ڈاک کے اہلکار} \textsuperscript{(44)} اور \VUgras{خط بانٹنے
والے} \textsuperscript{(45)} ڈاک خانے سے نکل رہے ہیں۔ ایک
\VUgras{ٹرام بان} \textsuperscript{(48)} کو اپنی گاڑی روکنی پڑی تاکہ
شہر کی \VUgras{ایمبولینس} \textsuperscript{(50)} گزر سکے، جو ایک
\VUgras{جرّاح} \textsuperscript{(52)} اور ایک \VUgras{تیمار دار}
\textsuperscript{(51)} کو لے کر آ رہی ہے تاکہ اُس \VUgras{زخمی}
\textsuperscript{(53)} کی مرہم پٹی کریں جسے \VUgras{اسٹریچر}
\textsuperscript{(54)} پر \VUgras{اخبار کے کھوکھے} \textsuperscript{(55)}
کے پاس لایا گیا ہے۔

%%K t11-c2-02-1 p
2. --- پیادہ فوج کا ایک دستہ، جو مشق سے واپس آ رہا تھا، رُک گیا تاکہ
\VUgras{شہری گھڑ سوار پولیس} \textsuperscript{(61)} کے ساتھ نظم و ضبط
قائم رکھے۔ \VUgras{گھڑ سوار}، جن کے گھوڑے اگلی ٹانگیں اٹھا رہے ہیں،
\VUgras{سپاہیوں} \textsuperscript{(57)} یا \VUgras{فوجی پولیس}
\textsuperscript{(63)} سے بہتر طور پر \VUgras{تماشائیوں}
\textsuperscript{(56)} کو پیچھے ہٹا رہے ہیں۔ ویسے بھی فوجی پولیس بہت
مصروف ہے۔ اُن میں سے ایک \VUgras{پولیس والوں} \textsuperscript{(64)}
کو بتا رہا ہے کہ ایک اور \VUgras{حادثے کے شکار} \textsuperscript{(65)}
کو کہاں لے جانا ہے — ایک بہادر آگ بجھانے والا، جو سیڑھی سے گر پڑا
ہے۔

%%K t11-c2-03-1 p
3. --- \VUgras{افسر} \textsuperscript{(66)} وہ جگہ بتا رہا ہے جہاں
آنے والا \VUgras{بھاپ کا پمپ} \textsuperscript{(67)} رکھنا ہے۔ اُس
زور دار مشین کا انتظار کرتے ہوئے اُس نے جلدی سے ایک \VUgras{ہاتھ کا
پمپ} \textsuperscript{(68)} لگوا دیا، جس کی لمبی \VUgras{نلی}
\textsuperscript{(69)} آگ پر ندی کی طرح پانی پھینکتی ہے۔

چھ آگ بجھانے والے اُسے چلا رہے ہیں، جب کہ اُن کا ساتھی، جو
\VUgras{نلی کا منہ} \textsuperscript{(70)} تھامے ہوئے ہے،
\VUgras{دھار} \textsuperscript{(71)} کو آگ کے بیچوں بیچ کی طرف موڑتا
ہے۔

%%K t11-c2-04-1 p
4. --- تماشا گھر کی دوسری طرف بڑی \VUgras{سیڑھی} \textsuperscript{(72)}
ٹکائی گئی ہے۔ اُس کے سہارے آگ بجھانے والے اُن شعلوں کے قریب پہنچ
سکتے ہیں جو کالے \VUgras{دھوئیں} \textsuperscript{(75)} کے بھنوروں کے
بیچ سے لپک رہے ہیں۔

%%K t11-tit-7 sub
\VUsaut{5.0mm}
\VUpk{10.2pt}{40}{بہادری کے کارنامے۔}
\VUsaut{3.4mm}

%%K t11-c2-05-1 p
5. --- \VUgras{آتش زدگی} \textsuperscript{(74)} \VUgras{تماشا گھر}
\textsuperscript{(73)} کے غلام گردش میں شروع ہوئی، اُس وقت جب
\VUgras{اوپیرا} کی مشق ہو رہی تھی، جس کا پہلا \VUgras{کھیل} کل ہونا
تھا۔ خوش قسمتی سے ہال میں \VUgras{ناظرین} \textsuperscript{(76)} تھوڑے
ہی تھے۔ بے چاروں نے \VUgras{چھجّے} \textsuperscript{(77)} پر پناہ لی،
جہاں بہادر \VUgras{آگ بجھانے والے} \textsuperscript{(86)} سیڑھی اور
موٹے رسّوں سے اُنھیں بچانے آتے ہیں۔

%%K t11-c2-06-1 p
6. --- \VUgras{ستون دار برآمدے} \textsuperscript{(78)} کے نیچے، جہاں
\VUgras{اشتہار} \textsuperscript{(79)} کھیل کی خبر دیتا ہے،
\VUgras{ساز نواز} \textsuperscript{(81)} بھاگتے دکھائی دیتے ہیں، جو
\VUgras{ارکسٹرا} \textsuperscript{(80)} کے ہیں ؛ وہ اپنے ساز اٹھائے
جا رہے ہیں : \VUgras{وائلن} \textsuperscript{(81)}، \VUgras{بانسری}
\textsuperscript{(82)}، \VUgras{چیلو} \textsuperscript{(83)}،
\VUgras{ٹرومبون} \textsuperscript{(84)}، \VUgras{ڈھول}
\textsuperscript{(85)} وغیرہ۔ لوگ \VUgras{ساز و سامان}
\textsuperscript{(87)} کا کچھ حصہ بچانے کی کوشش کر رہے ہیں۔

%%K t11-c2-07-1 p
7. --- \VUgras{اداکاروں} \textsuperscript{(89)} اور
\VUgras{اداکاراؤں} \textsuperscript{(88)} کو، \VUgras{گلوکاروں} اور
\VUgras{گلوکاراؤں} کو اپنے تماشے کے \VUgras{لباس}
\textsuperscript{(90)} ہی میں بھاگنا پڑا۔ ٹینر ایک \VUgras{صحافی}
\textsuperscript{(92)} کو بتا رہا ہے کہ یہ سب کیسے ہوا۔
\VUgras{نامہ نگار} پہلے ہی لمحے سے حکام کے ساتھ دوڑا آیا تھا :
\VUgras{ڈپٹی کمشنر} \textsuperscript{(91)}، \VUgras{ناظمِ شہر}
\textsuperscript{(93)}، \VUgras{پولیس کمشنر} \textsuperscript{(94)} اور
\VUgras{عدالتی افسر} \textsuperscript{(95)}، جو ذمہ داریوں کا فیصلہ
کرنے کے لیے تفتیش کریں گے۔

\VUsaut{6.0mm}
\VUfilet{22mm}
